\documentclass{exam-zh}

\examsetup{
  page = {
    size = a4paper,
    margin = 22mm,
  },
  choices = {
    margin = 8mm,
    columns = 2,
    column-sep = 2em,
    linesep = 0.8em,
  },
}

\definecolor{IssueRed}{HTML}{A63A32}
\definecolor{IssueRedBackground}{HTML}{FFF4F2}
\definecolor{IssueGreen}{HTML}{24734B}
\definecolor{IssueGreenBackground}{HTML}{F1F8F4}
\definecolor{IssueInk}{HTML}{202428}
\definecolor{IssueMuted}{HTML}{5E6872}
\definecolor{IssueRule}{HTML}{AEB7BF}

\newtcolorbox{issuepanel}[3]{
  enhanced,
  breakable=false,
  colback=#2,
  colframe=#1,
  boxrule=0.7pt,
  arc=1mm,
  left=5mm,
  right=5mm,
  top=4mm,
  bottom=4mm,
  before skip=4mm,
  after skip=5mm,
  title={#3},
  coltitle=white,
  fonttitle=\bfseries,
  colbacktitle=#1,
  boxed title style={arc=0.6mm},
}

\ExplSyntaxOn
\int_new:N \g__examzh_issue_forty_seven_row_int
\cs_new_eq:NN
  \__examzh_issue_forty_seven_apply_indent:
  \__examzh_choices_apply_left_indent:

\NewDocumentCommand \UseIssueFortySevenOldBehavior { }
  {
    \int_gzero:N \g__examzh_issue_forty_seven_row_int
    \cs_set_protected:Npn \__examzh_choices_apply_left_indent:
      {
        \int_gincr:N \g__examzh_issue_forty_seven_row_int
        \int_compare:nNnT
          { \g__examzh_issue_forty_seven_row_int } = { 1 }
          { \__examzh_issue_forty_seven_apply_indent: }
      }
  }
\ExplSyntaxOff

\newcommand\samechoicesexample{%
  \noindent\textbf{测试多行选项缩进：}
  \begin{choices}
    \item 第一行左侧选项
    \item 第一行右侧选项
    \item 第二行左侧选项
    \item 第二行右侧选项
  \end{choices}%
}

\begin{document}
\color{IssueInk}

\begin{center}
  {\LARGE\bfseries GitHub Issue \#47：修复前后对照}\par
  \vspace{1mm}
  {\large \texttt{choices/margin} 在多行选项中的缩进一致性}
\end{center}

\vspace{3mm}

\noindent
\textbf{问题描述。}
当 \texttt{choices} 使用两列、两行布局并设置
\texttt{margin=8mm} 时，旧实现只在环境开始处应用一次左缩进。
因此第一行 A/B 正常向右缩进，换行生成的第二行 C/D 却回到内容区左边界，
同一列的选项无法上下对齐。预期行为是每一行都使用相同的左右边距。

\begingroup
  \UseIssueFortySevenOldBehavior
  \begin{issuepanel}{IssueRed}{IssueRedBackground}
    {修复前：局部模拟 v0.3.4 的旧行为}
    \samechoicesexample
    \par\smallskip
    \noindent\color{IssueRed}\bfseries
    可见问题：C 比 A 更靠左，D 比 B 更靠左。
  \end{issuepanel}
\endgroup

\begin{issuepanel}{IssueGreen}{IssueGreenBackground}
  {修复后：当前工作区实现}
  \samechoicesexample
  \par\smallskip
  \noindent\color{IssueGreen}\bfseries
  修复结果：A/C 左边界一致，B/D 左边界一致。
\end{issuepanel}

\noindent
\textcolor{IssueMuted}{\textbf{测试条件：}}
\textcolor{IssueMuted}{XeLaTeX；\texttt{columns=2}；
\texttt{margin=8mm}；上下两部分使用完全相同的选项内容。}

\vfill
\noindent\textcolor{IssueRule}{\rule{\linewidth}{0.4pt}}
\par\smallskip
\noindent\footnotesize\color{IssueMuted}
上半部分只在局部分组中关闭后续行的缩进调用，用于复现问题；
下半部分未做模拟，直接使用当前 \texttt{exam-zh-choices.sty}。

\end{document}
